% Classical input: DLMF 25.12.12. Repository comparison in PROVENANCE.json.
\section{A local classification at reciprocal orders}
\label{sec:endpoints}

This section sharpens the nonintegral inverse chart and establishes
an obstruction to the extension proposed in Research Question~9 of
\emph{Resonant Order Jets and Convergent Inverse Transseries:
for Polylogarithms and Logarithmically Weighted Tails}
\cite{repo-inverse}. The starting expansion is classical; the
classification and uniform reciprocal-order chart below are deductions
from it. The source theorem is not contradicted: its global
endpoint-circle assertion is stated only for integer orders.

The nonintegral charts in the predecessor imply an exceptional phenomenon
that is worth separating from the integer-order resonance theory.
Write $s=1+\alpha>1$ and
\[
 D_\alpha(t)=\zeta(1+\alpha)-\operatorname{Li}_{1+\alpha}(e^{-t}).
\]
For nonintegral $\alpha$, the classical convergent expansion \cite{dlmf-polylog} gives
\[
 D_\alpha(t)=b(\alpha)t^\alpha+tA_\alpha(t),\qquad
 b(\alpha)=-\Gamma(-\alpha),\qquad
 A_\alpha(t)=\sum_{k\ge1}\frac{(-1)^{k+1}\zeta(1+\alpha-k)}{k!}t^{k-1}.
\]
Initially $0<|t|<2\pi$, $-\pi<\arg t<\pi$, and
$t^\alpha=\exp(\alpha\log t)$ uses the principal logarithm.
Further continuation is on the logarithmic covering.
The function $A_\alpha(t)$ is jointly holomorphic for $t$ near zero and
$\alpha$ in any sufficiently small neighborhood avoiding the nonnegative
integers.  Let $t_\alpha(y)$ denote the inverse which is positive for
positive small $y$.

\begin{theorem}[Minimal local ramification]\label{thm:ramification}
For real $s>1$ the minimal degree of a finite ramification $y=w^N$
making the endpoint inverse $t_{s-1}(y)$ holomorphic is as follows.
\begin{enumerate}
\item If $0<\alpha=s-1<1$ and $\alpha=p/q$ is in lowest terms,
the minimal degree is $p$.
\item If $\alpha>1$ is nonintegral and $\alpha=p/q$ is in lowest
terms, the minimal degree is $q$.
\item If $\alpha$ is irrational, or if $s$ is a positive integer
at least two, no finite ramification suffices.
\end{enumerate}
In particular, the endpoint inverse is holomorphic without ramification
if and only if $s=1+1/m$ for an integer $m\ge2$.
The same classification holds for the endpoint germ
$\operatorname{Li}_{s}^{\langle-1\rangle}(\zeta(s)-y)=e^{-t_{s-1}(y)}$.
\end{theorem}

\begin{proof}
Suppose first that $0<\alpha=p/q<1$.  The substitution $t=\xi^q$
makes the forward germ holomorphic:
\[
 D_\alpha(\xi^q)=b\xi^p+\xi^q A_\alpha(\xi^q).
\]
The factor multiplying $\xi^p$ is nonzero at zero, so an analytic
change of coordinate reduces this germ to $\eta^p$.  Its local inverse
is therefore holomorphic in $y^{1/p}$, and its $t$ coordinate starts
with $b^{-q/p}y^{q/p}$.  Since $p$ and $q$ are coprime, this leading
term forces every possible ramification degree to be divisible by $p$.
Thus the minimal degree is exactly $p$.

If $\alpha=p/q>1$ is nonintegral, the same substitution gives a
forward germ of order $q$, because $A_\alpha(0)=\zeta(\alpha)>0$.
It follows that the inverse $t$ is holomorphic in $y^{1/q}$.  To see
minimality, let $T(y)$ be the analytic inverse of $tA_\alpha(t)$.
Expansion of the inverse equation gives
\[
 t_\alpha(y)=T(y)-b\zeta(\alpha)^{-\alpha-1}y^\alpha+o(y^\alpha),
\]
after retaining the analytic terms of $T$ through the relevant degree.
The displayed fractional coefficient is nonzero, so $q$ divides every
possible ramification degree.

For irrational $\alpha<1$, the leading exponent $1/\alpha$ is
irrational.  For irrational $\alpha>1$, the preceding first
nonanalytic exponent is $\alpha$.  Neither can occur in a convergent
Puiseux germ.  At $s=2$,
$D_1(t)=t(1-\log t)+O(t^2)$, whence
$t_1(y)\sim y/\log(1/y)$, which is not Puiseux.
At integer $s=m+1\ge3$, the singular term of the forward deficit is
$(-1)^m t^m\log t/m!$, and its linear coefficient is $a=\zeta(m)>0$.
Therefore the inverse has a nonzero logarithmic term
$(-1)^{m+1}y^m\log y/(m!a^{m+1})$ in addition to its analytic part.
Finite ramification multiplies the logarithm by a nonzero integer
and does not remove it.
Finally $t\mapsto e^{-t}$ is biholomorphic near zero, so it preserves
these ramification degrees.
\end{proof}

\subsection{Explicit coefficients at the removable endpoints}

Fix an integer $m\ge2$, put $\alpha=1/m$ and $b_m=-\Gamma(-1/m)>0$,
and define
\[
 G_m(\xi)=b_m+\sum_{k\ge1}
 \frac{(-1)^{k+1}\zeta(1+1/m-k)}{k!}\xi^{mk-1}.
\]
Then $y=\xi G_m(\xi)$, $t=\xi^m$, and Lagrange inversion gives
the exact finite extraction formulas
\[
 [y^n]t_{1/m}(y)=\frac mn[\xi^{n-m}]G_m(\xi)^{-n},
\]
\[
 [y^n]e^{-t_{1/m}(y)}
 =-\frac mn[\xi^{n-m}]e^{-\xi^m}G_m(\xi)^{-n}\qquad(n\ge1).
\]
Coefficients with a negative extraction exponent vanish.  In particular
\[
 t_{1/m}(y)=b_m^{-m}y^m
 -m\zeta(1/m)b_m^{-2m}y^{2m-1}+O(y^{3m-2}).
\]
This last remainder follows because the next possible degree in
$t(y)$ is $3m-2$ (the term with $k=2$ first contributes at $3m-1$).
For the ordinary inverse endpoint germ, every derivative of order
$1,\ldots,m-1$ vanishes, and its $m$th derivative with respect to
$w=\zeta(s)-y$ equals
$(-1)^{m+1}m!b_m^{-m}$.

For example, at $s=3/2$ one has $b_2=2\sqrt\pi$ and
\[
 t_{1/2}(y)=\frac{y^2}{4\pi}
 -\frac{\zeta(1/2)}{8\pi^2}y^3
 +\frac{5\zeta(1/2)^2}{64\pi^3}y^4+O(y^5).
\]
Thus the square-root singularity of the forward polylogarithm has
become a quadratic stationary point of its holomorphic endpoint inverse.

\subsection{A uniform chart crossing a removable endpoint}

The apparent disappearance of branching is not uniform under order
perturbation.  Nevertheless a short exact equation gives a convergent
chart that is uniform at the reciprocal order.

\begin{samepage}
\begin{theorem}[Reciprocal-order crossover chart]\label{thm:reciprocal-chart}
Let $m\ge2$ and let $K$ be a compact real interval.  Put
\[
 \alpha=\frac1m+\epsilon,\quad T=\log(1/y),\quad
 \tau=\epsilon T,\quad h=T^{-1},\quad r=y^{m-1},\quad q=y^m.
\]
There is a function $U(h,r,q;\tau)$ holomorphic in its first three
variables near zero, with a complex neighborhood of $K$ for the last
variable and neighborhoods uniform on $K$, such that the physical
inverse satisfies
\[
 t_\alpha(y)=y^m U(T^{-1},y^{m-1},y^m;\epsilon T).
\]
Its value at the origin is
\[
 U(0,0,0;\tau)=b_m^{-m}e^{m^2\tau}.
\]
The Taylor series in the three independent small variables converges
after physical substitution for sufficiently small positive $y$.
\end{theorem}
\end{samepage}

\begin{proof}
Substitute $t=y^m U$ into $D_\alpha(t)=y$ and divide by $y$.
The exact equation is
\[
 b(1/m+h\tau)e^{-m\tau}U^{1/m+h\tau}
 +rU A_{1/m+h\tau}(qU)=1.
\]
For $h=r=q=0$ its indicated solution is
$U_0(\tau)=b_m^{-m}e^{m^2\tau}$.
Choose the logarithm
$\log U_0(\tau)=-m\log b_m+m^2\tau$ and write $U=U_0V$.
In this coordinate the equation is holomorphic near $V=1$, and its
$V$ derivative at $h=r=q=0$, $V=1$ is $1/m$.
The holomorphic implicit function theorem and a finite compactness
cover of $K$ give the claimed uniform chart.  Local uniqueness glues
the charts.  The physical positive branch is selected by continuity
from positive $y$ and the strictly increasing real forward deficit.
\end{proof}

For each fixed $j\ge1$, differentiation in $\alpha$ at fixed deficit
$y$ gives
\[
 \left.\partial_\alpha^j t_\alpha(y)\right|_{\alpha=1/m}
 =b_m^{-m}(-m^2)^j y^m(\log y)^j
 +O\bigl(y^m(1+|\log y|)^{j-1}\bigr).
\]
One may prove this either from the uniform chart or by differentiating
$b(\alpha)^{-1/\alpha}y^{1/\alpha}$ and the jointly holomorphic
two-generator remainder.  Derivatives of the latter acquire at most
fixed powers of $\log y$ and retain a positive additional power of
$y$.  Thus order differentiation exposes an infinite logarithmic
monodromy which is absent from the base germ.
Here the independent variable is the deficit $y$; differentiation at
fixed $w$ also differentiates $\zeta(s)-w$ and is a different operation.

\subsection{Why the integer endpoint-circle proof cannot extend verbatim}

For $0<\alpha=s-1<1/2$ and $u\downarrow0$, the upper-bank value is
\[
 \operatorname{Li}_{1+\alpha}(e^u+i0)
 =\zeta(1+\alpha)+\Gamma(-\alpha)e^{-i\pi\alpha}u^\alpha+O(u).
\]
Since $\Gamma(-\alpha)<0$ and $\cos(\pi\alpha)>0$, this yields
\[
 |\operatorname{Li}_{s}(e^u+i0)|
 =\zeta(s)+\Gamma(1-s)\cos(\pi(s-1))u^{s-1}
 +o(u^{s-1})<\zeta(s)
\]
for every sufficiently small positive $u$.
Consequently the cut-modulus assertion in the integer-order theorem
fails for every $1<s<3/2$.  The predecessor states that assertion only
for integers, so this is a negative answer to a proposed extension,
not a counterexample to a printed theorem.

The cut entering the endpoint circle does not itself determine the
Taylor radius of an analytically continued inverse: the cut is a
boundary of the chosen forward sheet, not necessarily an inverse
singularity.  Likewise endpoint removability does not rule out another
singularity of the same modulus.  Determining the global radius at the
reciprocal orders remains a separate problem.
